\documentclass[11pt]{article}
\usepackage{amsmath,amssymb,amsthm,mathtools,hyperref}
\usepackage[margin=26mm]{geometry}
\newtheorem{theorem}{Theorem}
\newtheorem{proposition}[theorem]{Proposition}
\newtheorem{lemma}[theorem]{Lemma}
\DeclareMathOperator{\Gram}{Gram}
\DeclareMathOperator{\Res}{Res}
\DeclareMathOperator{\im}{im}
\DeclareMathOperator{\cond}{cond}
\title{Original endpoint retractions and the reflected-pair period calculation}
\author{Split-Zero research programme}
\date{19 September 2026}
\begin{document}
\maketitle

\section{Objects, source conventions, and the receiving maps}
Keep a complete stipulated quartet
\[
 \rho=\tfrac12+\delta+i\gamma,\quad
 \mathcal R=\{\rho,\bar\rho,1-\rho,1-\bar\rho\},\quad
 0<\delta<\tfrac12,\quad\gamma>2,
\]
each of order \(m\). Set
\[
 h(s)=\prod_{\alpha\in\mathcal R}(s-\alpha)^m,\quad
 v_h=2\xi/h,\quad U_h=j_hv_h,\quad
 E_h=\mathbb C[s]/(h).
\]
For \(k\equiv1\pmod4,\ k\ge5\), retain
\[
 e=1+k(m-1),\quad c=k/2,\quad q=e(k+1)^2,\qquad
 \chi(S)=\prod_{a,b=0}^k
 \bigl(S-c-(2a-k)\delta-i(2b-k)\gamma\bigr)^e.
\]
The spaces and operators are \(E_k=\mathbb C[S]/(\chi)\),
\(A=M_S\), \(A_\Sigma=\sum_iM_{s_i}\), and
\[
 \eta_k:E_k\longrightarrow E_h^{\otimes k},\qquad
 [P]\longmapsto U_h^{\otimes k}P(A_\Sigma)1.                 \tag{ET1}
\]
These are the original maps OPG1--5, not altered unit values
\cite[OPG1--5]{Current}. The exact source is
\[
 w_h(y)=|v_h(1/2+iy)|^2/(2\pi),\qquad m_k=w_h^{*k},\qquad
 \sigma(y)=|\Gamma(1/4+iy/2)|^2/(2\pi).
                                                               \tag{ET2}
\]
All norms below keep their stated measures. In particular, a coordinate
congruence used to calculate eigenvalues does not change either metric.

\section{The complete finite endpoint retraction}
Write \(\zeta(\alpha+z)=z^m z_\alpha(z)\), \(z_\alpha(0)\ne0\).
The finite singular germ attached to a local class \(u\) is exactly
\[
 \mathfrak b_\alpha u
 =x^{-\alpha}\Res_{z=0}
       \frac{e^{zL}u(z)}{2z^mz_\alpha(z)}\,dz
 =x^{-\alpha}\sum_{j=0}^{m-1}
 [z^{m-1-j}]\frac{u(z)}{2z_\alpha(z)}\frac{L^j}{j!},
 \quad L=\log(1/x).                                    \tag{ET3}
\]
The equality follows by multiplying the convergent Taylor germs and
selecting the coefficient of \(z^{-1}\). Thus no unknown lower local
coefficient has been suppressed. This verifies the finite-germ formula
E1 of the incoming note \cite[E1]{Incoming}. Identifying it with a
particular infinite-dimensional analytic connecting map additionally
uses that map's specified source extension; the construction here works
directly with the displayed residue and does not change its domain.

The linear reversal \(J_m(z^r)=L^{m-1-r}/(m-1-r)!\) satisfies
\(J_mM_z=\partial_LJ_m\), including the top nilpotent power. Multiplication
by the entire unit commutes with \(M_z\). Hence ET3 intertwines local
multiplication by \(\alpha+z\) with \(\alpha+\partial_L\).
For an ordered tuple \(\boldsymbol\alpha\), put
\(\beta=\sum_i\alpha_i\), \(n=k(m-1)\). Diagonal multiplication is
\[
 \mathfrak d_{\boldsymbol\alpha}
 \left(\bigotimes_i x_i^{-\alpha_i}L_i^{j_i}/j_i!\right)
 =x^{-\beta}\frac{(\sum_i j_i)!}{\prod_i j_i!}
                  \frac{L^{\sum_i j_i}}{(\sum_i j_i)!}.       \tag{ET4}
\]
The product rule proves exact intertwining with
\(\sum_i(\alpha_i+\partial_{L_i})\).
In particular the factor \(n!/((m-1)!)^k\) on the longest polynomial
does not disappear.

Let \(A_*(s)=s(s-1)\pi^{-s/2}\Gamma(s/2)/2\) and
\(h_\alpha(s)=h(s)/(s-\alpha)^m\).
The identity \(\xi=A_*\zeta\) gives, in the original local algebra,
\[
 \frac{U_h(\alpha+z)}{2z_\alpha(z)}
       =\frac{A_*(\alpha+z)}{h_\alpha(\alpha+z)}
                  \pmod {z^m}.                            \tag{ET5}
\]
All these units are invertible: the stipulated roots are nontrivial,
the Gamma function has no zeros and no pole at these points
\cite[\S5.2(i)]{DLMF}, and
the full orders have already been divided out.
Define the actual polynomial
\[
 C_{\boldsymbol\alpha}(L)=
 \prod_{i=1}^k\left(
  \sum_{j=0}^{m-1}[z^{m-1-j}]
  \frac{A_*(\alpha_i+z)}{h_{\alpha_i}(\alpha_i+z)}
                  \frac{L^j}{j!}\right)
 =\sum_{j=0}^n c_jL^j/j!.
                                                               \tag{ET6}
\]
Its top coefficient is
\[
 c_n=\frac{n!}{((m-1)!)^k}
                 \prod_i\frac{U_h(\alpha_i)}{2z_{\alpha_i}(0)}
 \ne0.                                                        \tag{ET7}
\]
With \(D_{\boldsymbol\alpha}=\mathfrak d_{\boldsymbol\alpha}
\bigotimes_i\mathfrak b_{\alpha_i}\),
\[
 D_{\boldsymbol\alpha}\eta_kP
       =x^{-\beta}P(\beta+\partial_L)C_{\boldsymbol\alpha}(L).
                                                               \tag{ET8}
\]
For an arbitrary output \(x^{-\beta}\sum b_jL^j/j!\), set
\(B(T)=\sum_{r=0}^n b_{n-r}T^r\) and
\(H(T)=\sum_{r=0}^n c_{n-r}T^r\). Reversal intertwines \(\partial_L\)
with multiplication by \(T\), so ET8 is precisely
\[
 B(T)=P(\beta+T)H(T)\pmod{T^{n+1}}.
\]
The complete inverse is finite:
\[
 P(\beta+T)=H(T)^{-1}B(T),\qquad
 H(T)^{-1}=c_n^{-1}\sum_{j=0}^n
                  \left(-\frac{H(T)-c_n}{c_n}\right)^j
                   \pmod{T^{n+1}}.                         \tag{ET9}
\]

Choose and record one ordered tuple at every original centre
\(\beta_{ab}=c+(2a-k)\delta+i(2b-k)\gamma\).
Every such tuple exists because the two signs in a quartet root are
independent. Restrict to that tuple, apply ET9, and use the ordinary
Chinese remainder inverse for the pairwise coprime polynomials
\((S-\beta_{ab})^e\). This defines the exact map
\[
 R_\partial:E_h^{\otimes k}\longrightarrow E_k,\qquad
 R_\partial\eta_k=I,\qquad R_\partial A_\Sigma=AR_\partial.
                                                               \tag{ET10}
\]
These two identities follow component by component from ET9 and from
the derivative intertwining. They also prove injectivity of ET1
without a dimension argument. All unselected tuples remain in the
domain and belong to the specified kernel. For any earlier
equivariant retraction \(R\) of this same inclusion, the exact
comparison is
\[
 R_\partial-R=(D\eta_k)^{-1}D(I-\eta_kR),                 \tag{ET11}
\]
where \(D\) is the direct sum of the selected maps, with the
Chinese remainder target understood. Indeed its right side is
\(R_\partial-R_\partial\eta_kR=R_\partial-R\).
This is an explicit map through the old complementary module;
neither complement is asserted orthogonal.

The complete Jordan list of that kernel can be recovered directly.
Put \(d=m-1\) and give the auxiliary polynomial basis \(L^j\),
\(0\le j\le d\), squared lengths \(1/\binom dj\). This auxiliary
inner product is used only to prove a finite representation identity;
it is not substituted for ET2. The operators
\[
 N=\partial_L,\quad H=2L\partial_L-d,\quad
 R=L(d-L\partial_L)
\]
satisfy \(N^*=R\), \(H^*=H\),
\([H,N]=-2N\), \([H,R]=2R\), \([R,N]=H\).
Their sums act on the ordered tensor. If \(Rv=0\), \(Hv=\ell v\),
then
\[
 RN^jv=j(\ell-j+1)N^{j-1}v,\qquad
 \|N^jv\|^2=j(\ell-j+1)\|N^{j-1}v\|^2.
\]
These identities follow by induction from the commutators and
adjointness. Taking a highest nonzero weight in any invariant
orthogonal complement therefore constructs a string of length
\(\ell+1\); its orthogonal complement is again invariant because
both adjoints belong to the operator family. Induction on the finite
dimension decomposes the complete tensor into these strings.
If \(p_r=[t^r](1+\cdots+t^d)^k\), the weight multiplicities and their
symmetry give exactly \(p_r-p_{r-1}\) strings of length \(kd-2r+1\),
for \(0\le r\le\lfloor kd/2\rfloor\), with \(p_{-1}=0\).
In particular the longest string is unique. Diagonal multiplication
ET4 intertwines all three displayed operators with those of degree
\(kd\), by the product rule. It is onto by its nonzero top value,
and kills each shorter string: a highest weight vector of smaller
weight maps to a vector killed by \(R\) in that same smaller weight,
where the degree-\(kd\) target has no such vector.

At \(\beta_{ab}\), there are exactly
\(M_{ab}=\binom ka\binom kb\) ordered tuples, by choosing the two
sign sets independently. ET10 removes one longest string from their
direct sum. Its kernel consequently has \(M_{ab}-1\) strings of
length \(e\) and
\[
 M_{ab}(p_r-p_{r-1})\quad\hbox{strings of length }e-2r
 \quad(1\le r\le\lfloor(e-1)/2\rfloor).              \tag{ET11a}
\]
Multiplication by the full local units commutes with the original
nilpotent action, so it changes the explicit splitting but preserves
this complete list.

\section{Actual period coordinates and single-class survival}
This section uses the original simple lane \(m=1\). Put
\[
 x=\delta+i\gamma,\quad
 h_c(w)=w^4+\beta w^2+\eta,\quad
 \beta=2(\gamma^2-\delta^2),\quad\eta=(\delta^2+\gamma^2)^2,
\]
\[
 e_x(w)=\frac{h_c(w)}{(w-x)h_c'(x)}
 =\frac{w^3+xw^2+(x^2+\beta)w+x(x^2+\beta)}{h_c'(x)}.
                                                               \tag{ET12}
\]
Use its column \(\mathbf v_x\) in the unchanged frame
\(1,w,w^2,w^3\). The recovered provider gives
\[
 P(u)=e^{\mathfrak a/u}G_u\mathcal R(u^{-1}),\quad
 G_u=\operatorname{diag}(u^{a/5}g_a)_{a=1}^4,\quad
 g_a=5^{a/5-1}e^{\pi ia/5}\Gamma(a/5),
                                                               \tag{ET13}
\]
\[
 \mathfrak a=1/160+\beta/24+\eta/2,\quad
 \det\mathcal R(z)=1,\quad
 \mathcal R(0)=
 \begin{pmatrix}1&0&-\beta/3&0\\0&1&0&-2\beta/3\\
 0&0&1&0\\0&0&0&1\end{pmatrix}=:J_\infty .
                                                               \tag{ET14}
\]
The exact coefficients and their convergent bound are PCL2--6
\cite[PCL1--6, RPC8]{Current}:
\[
 R_{ar,n}=
 \sum_{\substack{p,h\ge0\\2p+4h=5n+r+1-a}}
 \frac{\beta^p\eta^h}{3^pp!h!}(-1)^{p+h-n}
 5^{p+h-n}(a/5)_{p+h-n},\quad
 C_1=\sum_{n\ge1}\|R_n\|_1<\infty ,
                                                               \tag{ET15}
\]
\[
 \|\mathcal R(z)-J_\infty\|_1\le C_1|z|\quad(|z|\le1).
\]
The original logarithm branch remains in \(G_u\). Let
\(\zeta_5=e^{2\pi i/5}\), \(D=\operatorname{diag}(\zeta_5^a)\) and
\[
 \mathcal P_k=\frac15\sum_{j=0}^4(D^{-j})^{\otimes k}.
                                                               \tag{ET16}
\]
OPG5 and its Fourier conjugacy show that this is the actual
observation projector in the stated Fourier coordinates, not a new
observation \cite[OPG5, PQO5]{Current}. It is an orthogonal projector
for the Euclidean metric in those coordinates.
Writing \(b=P(u)\mathbf v_x\), the original extremal class has output
\[
 U_h(\rho)^k\mathcal P_kb^{\otimes k}.                  \tag{ET17}
\]
Since \(P(u)\) is invertible, \(b\ne0\). For \(5\mid k\), choose
any nonzero entry \(b_a\). The all-\(a\) ordered tensor entry is
nonzero and survives ET16. Thus ET17 is nonzero on the original
cofinal sequence \(k=20n+5\). If two entries \(b_a,b_b\) are nonzero,
choose \(j\in\{0,\ldots,4\}\) with \(ka+j(b-a)=0\bmod5\).
For \(k\ge4\), the entry with \(j\) labels \(b\) and \(k-j\) labels
\(a\) exists and proves survival.

With \(p_a=|b_a|^2/\|b\|^2\), exact character orthogonality gives
\[
 f_k(b):=\frac{\|\mathcal P_kb^{\otimes k}\|^2}{\|b\|^{2k}}
  =\frac15\sum_{j=0}^4
       \left(\sum_a p_a\zeta_5^{ja}\right)^k.          \tag{ET18}
\]
If at least two entries are nonzero,
\(r_b=\max_{1\le j\le4}|\sum_ap_a\zeta_5^{ja}|<1\),
since distinct fifth roots cannot all have the same argument.
Consequently \(|f_k(b)-1/5|\le4r_b^k/5\).
If exactly one entry is nonzero, \(f_k(b)\) is one for \(5\mid k\)
and zero otherwise. These statements verify E15--18 in their
actual period metric, with the common unit cancelling only in the ratio.

\section{A new reflected-pair observation calculation}
Retain \(b_+=P(u)\mathbf v_x\), \(b_-=P(u)\mathbf v_{-x}\) and
\[
 a=\|b_+\|^2,\quad b_0=\|b_-\|^2,\quad
 z_j=b_+^*D^jb_-,\quad
 a_j=b_+^*D^jb_+,\quad b_j=b_-^*D^jb_- .
\]
The exact output Gram on the original reflected pair is
\[
 \mathcal O_k=
 |U_h(\rho)|^{2k}\frac15
 \begin{pmatrix}
 \sum_{j=0}^4a_j^k&\sum_{j=0}^4z_j^k\\
 \sum_{j=0}^4\overline{z_j}^{\,k}&\sum_{j=0}^4b_j^k
 \end{pmatrix}.                                      \tag{ET19}
\]
The lower left uses the summed conjugates; replacing \(j\) by
\(-j\) proves Hermitian symmetry directly. To prove ET19, insert
ET16 between each pair of tensor vectors and use
\(\langle v^{\otimes k},w^{\otimes k}\rangle
=\langle v,w\rangle^k\).
The original reflection has the same unit at \(\rho\) and \(1-\rho\)
because both \(h\) and \(\xi\) satisfy the same reflection; this is
why the displayed unit is common. No cross term is discarded.

Here is an explicit nonalignment guard derived from the actual PCL
coefficients. Put
\[
 A_0=x/h_c'(x),\quad B_0=1/h_c'(x),\quad
 c_5=\min_j|1+\zeta_5^j|=(\sqrt5-1)/2,
\]
\[
 V=\|\mathbf v_x\|_1=\|\mathbf v_{-x}\|_1,\qquad
 U_*=\max\left\{1,\frac{16C_1V}
                    {c_5\min\{|A_0|,|B_0|\}}\right\}. \tag{ET20}
\]
The last two entries of \(J_\infty\mathbf v_x\) are
\((A_0,B_0)\), and those of \(J_\infty\mathbf v_{-x}\) are
\((A_0,-B_0)\), because \(h_c'(-x)=-h_c'(x)\).
For \(|u|\ge U_*\), every perturbation of those four entries is at
most \(d=c_5\min(|A_0|,|B_0|)/16\).
For each \(j\), the two-row minor between these two columns after
multiplying the second by \(D^j\) has unperturbed modulus
\(|A_0B_0||1+\zeta_5^j|\).
Its perturbation is at most
\[
 2d(|A_0|+|B_0|)+2d^2
 \le |A_0B_0|(c_5/4+c_5^2/128)
 <\tfrac12c_5|A_0B_0|.
                                                               \tag{ET21}
\]
Thus all five minors are nonzero. Multiplication by the original
nonzero diagonal \(e^{\mathfrak a/u}G_u\) does not change their
nonvanishing. In particular \(b_+\) and \(D^jb_-\) are independent
for every \(j\). This also guarantees two nonzero entries in both
columns.

At a fixed admitted period satisfying ET20 define the actual numbers
\[
 r_+=\max_{1\le j\le4}|a_j|/a<1,\quad
 r_-=\max_{1\le j\le4}|b_j|/b_0<1,\quad
 r_\times=\max_{0\le j\le4}|z_j|/\sqrt{ab_0}<1.
                                                               \tag{ET22}
\]
Strictness in the third inequality follows from ET21 and the
equality case of Cauchy--Schwarz. These are finite sums of the four
original coordinates; no unknown asymptotic Gram replaces them.
On the finite guards \(4r_+^k\le1\) and \(4r_-^k\le1\),
ET19 gives the lower determinant bound
\[
 \frac{\det\mathcal O_k}
 {|U_h(\rho)|^{4k}(ab_0)^k}
 \ge\frac{(1-4r_+^k)(1-4r_-^k)}{25}
                 -r_\times^{2k}.                    \tag{ET23}
\]
For example the right side is positive whenever
\[
 4r_+^k\le\tfrac14,\qquad4r_-^k\le\tfrac14,\qquad
 25r_\times^{2k}<\tfrac9{16}.                         \tag{ET24}
\]
Taking logarithms supplies a finite threshold directly from ET22;
if a radius is zero its condition is already satisfied.
This proves eventual injectivity of the original observation on the
entire reflected pair at every fixed admitted period in ET20.
It strengthens single-class nonvanishing without asserting a
characteristic-changing comparison.

There is a quantitative full-pair version of the \(1/5\) calculation.
The unprojected period Gram is
\[
 \mathcal H_k=|U_h(\rho)|^{2k}
       \begin{pmatrix}a^k&z_0^k\\\bar z_0^{\,k}&b_0^k\end{pmatrix}.
\]
Define
\[
 \epsilon_k=
 \frac{\frac45\max(r_+^k,r_-^k)+\frac45r_\times^k}
 {1-r_\times^k}.
                                                               \tag{ET25}
\]
Then
\[
 (1/5-\epsilon_k)\mathcal H_k\preceq
       \mathcal O_k\preceq(1/5+\epsilon_k)\mathcal H_k.
                                                               \tag{ET26}
\]
Indeed congruence by
\(|U_h(\rho)|^{-k}\operatorname{diag}(a^{-k/2},b_0^{-k/2})\)
leaves the unprojected Gram with diagonal one and off-diagonal
modulus at most \(r_\times^k\).
The difference \(\mathcal O_k-\mathcal H_k/5\) has diagonal entries
at most \(4\max(r_+^k,r_-^k)/5\) in modulus, and off-diagonal
entry at most \(4r_\times^k/5\), since its \(j=0\) summand cancels
exactly. Its operator norm is bounded by their sum. Dividing by
the lower eigenvalue \(1-r_\times^k\) proves ET26, after reversing
the congruence. Every unit and coordinate factor is therefore
present in the asserted original inequality.

\section{Reconstructing the first-cutoff arithmetic parity coefficient}
This proof recovers S2's parity coefficient from the source envelopes,
without using an unavailable separate parity calculation.
Set \(\lambda=k\rho=c+\mu\), \(\mu=k(\delta+i\gamma)\),
\(R=|\mu|\), and keep the terminal class
\[
 v=\left[\frac{\chi(S)}
  {(S-\lambda)\chi_\lambda(\lambda)}\right],\qquad
 \chi_\lambda=\chi/(S-\lambda)^e .
\]
At \(N=q-1\) the quotient representative is unique. Define the
complete polynomial
\[
 B_k(y)=\frac{\chi(c+iy)}{(iy-\mu)(iy+\mu)},\qquad n=q-2,\qquad
 M_j=\int_{\mathbb R}y^j|B_k(y)|^2m_k(y)\,dy .
                                                               \tag{ET27}
\]
Its leading coefficient has modulus one; its remaining \(n\) roots,
with their full remaining multiplicities, have modulus at most
\(R=k\sqrt{\delta^2+\gamma^2}\).
The polynomial is even. The literal representative is
\((iy+\mu)B_k(y)/\chi_\lambda(\lambda)\).
Evenness of the source kills the even/odd cross term exactly.
Consequently the even-energy fraction is
\[
 \theta_{q-1}=\frac{R^2M_0}{M_2+R^2M_0}.              \tag{ET28}
\]
This derivation is valid for every complete multiplicity; it never
uses \(\chi'(\lambda)\).

For completeness the source estimates used here are exactly GTM1
and RTM2 \cite[GTM1, RTM2]{Envelopes}:
\[
 a_k(1+y^2)^{-M}\sigma(y)\le m_k(y)\le A_k\sigma(y),
 \quad M=21+4m,\quad
 \log(A_k/a_k)=O_h(k+\log k),
                                                               \tag{ET29}
\]
\[
 c_\Gamma e^{-\alpha|y|}(1+|y|)^{-1/2}
 \le\sigma(y)\le C_\Gamma e^{-\alpha|y|}(1+|y|)^{-1/2},
\quad\alpha=\pi/2,
\]
\[
 c_\Gamma=\sqrt2e^{-7/6},\qquad
 C_\Gamma=\sqrt{2\sqrt5}e^{2/3}.
\]
In particular the full source mass remains in \(a_k,A_k\).
Their definitions are
\[
 a_k=\frac{c_h\vartheta_h^{k-3}e^{-\alpha(k-3)}(k-2)^{-42-8m}}
 {C_\Gamma 2^{21+4m}},\quad
 A_k=\frac{C_h^{\rm tilt}}{c_\Gamma}k^{8m-7/2}
                    \mathfrak M^{k-1},
\]
\[
 \vartheta_h=\int_{-1}^1w_h(y)\,dy,\qquad
 \mathfrak M=\int_{\mathbb R}e^{\alpha y}w_h(y)\,dy.
\]

Here are explicit finite bounds establishing concentration, including
the unbounded tails. Set \(x_0=2n/\alpha\), and impose only the
eventual guard \(R<x_0\). Restriction to \([x_0,x_0+1]\) gives
\[
 M_0\ge L_k:=
 \frac{a_kc_\Gamma e^{-\alpha(x_0+1)}(x_0-R)^{2n}}
 {(1+(x_0+1)^2)^M\sqrt{x_0+2}}.                      \tag{ET30}
\]
For \(j=0,2\), the upper bound on either half-line is
\[
 |y|^j|B_k(y)|^2m_k(y)
 \le A_kC_\Gamma |y|^j(|y|+R)^{2n}e^{-\alpha|y|}.
                                                               \tag{ET31}
\]
Put \(s_*=2/\alpha=4/\pi\) and fix \(0<\varepsilon<s_*\).
For \(t=|y|+R\), the region
\(\bigl||y|/q-s_*\bigr|>\varepsilon\) lies in
\[
 t<q(s_*-\varepsilon)+R
 \quad\hbox{or}\quad t>q(s_*+\varepsilon)+R.
\]
For \(\nu_j=2n+j+1\), let \(T_j\) denote the auxiliary integral
with density \(\alpha^{\nu_j}t^{\nu_j-1}e^{-\alpha t}/(\nu_j-1)!\)
on \(t>0\); this auxiliary notation does not rescale ET2.
Integrating ET31 yields the exact bound
\[
 \frac{\int_{\{||y|/q-s_*|>\varepsilon\}}
              |y|^j|B_k(y)|^2m_k(y)\,dy}{q^jM_0}
 \le
 \frac{2A_kC_\Gamma e^{\alpha R}(2n+j)!}
 {q^j\alpha^{2n+j+1}L_k}\,
 \Pr\{T_j\notin[q(s_*-\varepsilon)+R,
                         q(s_*+\varepsilon)+R]\}.       \tag{ET32}
\]
For \(v>1\) and \(0<v<1\), respectively, elementary exponential
Markov estimates give
\[
 \Pr\{T_j\ge\nu_jv/\alpha\}\le
 e^{-\nu_j(v-1-\log v)},\qquad
 \Pr\{T_j\le\nu_jv/\alpha\}\le
 e^{-\nu_j(v-1-\log v)}.                             \tag{ET33}
\]
To prove them integrate \(e^{s\alpha t}\) against the displayed
density, obtaining \((1-s)^{-\nu_j}\), and choose \(s=1-1/v\).
This is positive for the upper tail and negative for the lower tail.
All integrals follow by repeated integration by parts.

The logarithm of the prefactor in ET32 is \(O_h(k+\log q)\).
Indeed \(\log(2n+j)!=(2n+j)\log(2n+j)-(2n+j)+O(\log n)\),
obtained by bounding the sum of \(\log r\) between its two adjacent
integrals; \(R=O_h(k)\), \(R/n\to0\), and
\(-2n\log(1-R/x_0)=O_h(k)\).
The remaining factors in ET30 are polynomial in \(q\), and ET29
gives the stated source ratio. Each tail in ET33 has exponent
\(-c_\varepsilon q+o(q)\), since its \(v\) tends to
\(1\pm\alpha\varepsilon/2\). Thus the right side of ET32 tends to
zero for \(j=0\) and \(j=2\). On the complementary region
\((s_*-\varepsilon)^2\le y^2/q^2\le(s_*+\varepsilon)^2\).
Letting \(\varepsilon\downarrow0\) proves
\[
 \frac{M_2}{q^2M_0}\longrightarrow\frac{16}{\pi^2},\qquad
 \theta_{q-1}=
 \frac{\pi^2(\delta^2+\gamma^2)}{16}
                  \frac{k^2}{q^2}(1+o_h(1)).         \tag{ET34}
\]
This proof used the unchanged complete root polynomial, all roots in
the degree count, both half-lines, and the actual source envelopes.

\section{Exact arithmetic pair metrics and their period comparison}
Reflection \(\Sigma P(S)=P(k-S)\) is an isometry of the original
quotient minimum because ET2 is even and \(\chi(k-S)=\chi(S)\).
For \(P_N\), its minimum representative of \(v\), put
\(E_N=\|P_N\|^2\) and
\(\theta_N=\|(P_N+\Sigma P_N)/2\|^2/E_N\).
The reflected pair is independent, and its exact Gram is
\[
 G_N^{\rm pair}=E_N
  \begin{pmatrix}1&h_N\\h_N&1\end{pmatrix},
 \qquad h_N=2\theta_N-1,\quad0<\theta_N<1.             \tag{ET35}
\]
The decomposition into even and odd parts proves the off-diagonal
entry and positivity. On this pair \(A-cI=\operatorname{diag}(\mu,-\mu)\).
The spectral coordinate projector onto the first line has norm
\[
 \|E_+\|_{G_N^{\rm pair}}=(1-h_N^2)^{-1/2}.           \tag{ET36}
\]
For a fixed first coordinate \(z_1\), minimization over \(z_2\) in
the denominator gives \((1-h_N^2)|z_1|^2\), proving the formula.
For real \(\tau\), \(T=\exp(\tau(A-cI))\) has squared singular
values with product one and sum
\[
 t_N(\tau)=\frac{2\cosh(2k\delta\tau)
                     -2h_N^2\cos(2k\gamma\tau)}{1-h_N^2}.
                                                               \tag{ET37}
\]
This follows by multiplying \(G^{-1}T^*GT\); therefore its two
eigenvalues are \((t_N\pm\sqrt{t_N^2-4})/2\).
Using ET34, for each fixed \(\tau\ne0\),
\[
 \|T\|_{G_{q-1}^{\rm pair}}
 =\frac{2q}{\pi k\sqrt{\delta^2+\gamma^2}}
              e^{|\tau|k\delta}(1+o_{h,\tau}(1)).     \tag{ET38}
\]
The original uncentred action has the additional factor
\(e^{\tau k/2}\), retained exactly.

For comparison with the proposed normal-action metric in S9, any
positive metric making this two-eigenvalue action normal is diagonal
in the displayed eigenbasis: \(H=\operatorname{diag}(a_*,b_*)\),
\(a_*,b_*>0\). Distinct eigenvectors of a normal operator are
orthogonal, as follows by applying the adjoint eigenvalue equation
to their inner product; conversely a diagonal metric makes this
diagonal action normal. The squared generalized trace divided by
the generalized determinant is
\[
 \frac{(a_*+b_*)^2}{a_*b_*(1-h_N^2)}.
\]
Its minimum is at \(a_*=b_*\), since
\((a_*-b_*)^2\ge0\). The positive generalized eigenvalue ratio
increases with that expression. Thus, exactly,
\[
 \inf_{H:\,A\ {\rm normal\ in}\ H}
 \cond((G_N^{\rm pair})^{-1}H)
 =\frac{1+|h_N|}{1-|h_N|},
 \quad
 \inf_{H:\,A\ {\rm normal\ in}\ H}
 \cond((G_{q-1}^{\rm pair})^{-1}H)
 =\frac{16q^2}{\pi^2(\delta^2+\gamma^2)k^2}
                      (1+o_h(1)).                    \tag{ET38a}
\]
This evaluates the metric cost of that specified construction; the
calculation of the original arithmetic metric still uses ET35.

In the simple period lane the unprojected period metric on the pair
is \(\mathcal H_k\) from ET25. Its relative generalized eigenvalue
condition number \(\kappa\) obeys the exact equation
\[
 \kappa+\kappa^{-1}+2=
 \frac{(a^k+b_0^k-2h_N\Re z_0^k)^2}
 {(1-h_N^2)((ab_0)^k-|z_0|^{2k})}.                  \tag{ET39}
\]
Indeed the right side is the squared generalized trace divided by
the generalized determinant. Both scalar factors
\(E_N\) and \(|U_h(\rho)|^{2k}\) cancel only in this condition number.
At every fixed admitted period, independence of \(b_+,b_-\) gives
\(|z_0|/\sqrt{ab_0}<1\). Combining ET34 with ET39 gives
\[
 \kappa_{q-1,k}=
 \frac{4q^2}{\pi^2(\delta^2+\gamma^2)k^2}
 \left[(a/b_0)^k+(b_0/a)^k+2\right](1+o_{h,u}(1)).
                                                               \tag{ET40}
\]
For example, the cross term in the trace divided by
\(a^k+b_0^k\) is at most
\((|z_0|/\sqrt{ab_0})^k\); the determinant correction is its
square, and the positive solution of ET39 divided by its right
side tends to one. These estimates prove ET40 also when the period
norms are unbalanced. Hence S14 is recovered from explicit source
envelopes and the original period provider, rather than imported
as an unsupported asymptotic.

\section{The exact next receiving calculation}
Let \(W=P(u)^{\otimes k}\eta_kE_k\) be the original cyclic image
in Fourier tensor coordinates and \(L=\mathcal P_k|_W\).
The period-metric quotient norm of the original output \(Lv\) is
the minimum of \(\|w\|^2\) over \(w\in W\) with \(Lw=Lv\).
As \(\mathcal P_k\) is orthogonal, it is at least \(\|Lv\|^2\)
and at most \(\|v\|^2\); this proves precisely the bounds in E18.
ET19--26 calculate \(L^*L\) on the reflected pair, including its
nonzero cross term. They do not replace the full cyclic fibre by
that two-dimensional domain.

In any fixed original basis of \(E_k\), with its full period Gram
\(H\) and the original observation matrix \(\Lambda\), the attained
quotient matrix in a full basis of the actual image is
\[
 Q_{\rm per}=(\Lambda H^{-1}\Lambda^*)^{-1}.           \tag{ET41}
\]
To prove this, set
\(S=H^{-1}\Lambda^*(\Lambda H^{-1}\Lambda^*)^{-1}\).
Then \(\Lambda S=I\) and \(S^*HI=0\) for every kernel column \(I\),
so \(Sv\) is the unique minimum lift and \(S^*HS=Q_{\rm per}\).
The arithmetic receiver has the identical exact formula with the
original full arithmetic \(G_N\). This specifies the two exact
maps whose actual inverse Grams must enter the seven allocations;
ET40 is the proven pair restriction of their metric comparison.
No individual arithmetic projected sign follows merely from ET18
or ET26, and none is assigned here.

The incoming E19--21 additionally refer to a finite-field exponential
family and a lower-boundary exact sequence. Deligne's weight theorem
is a theorem about its specified finite-field complexes and Frobenius
\cite[Thm.\ 3.3.1]{Deligne}. ET3--41 do not change coefficient fields
or assign Frobenius to the analytic dilation. The exact identity
\(\pi C\partial^{\otimes k}=\pi(C\partial^{\otimes k}-\iota b)\)
uses \(\pi\iota=0\); it neither constructs the comparison \(C\) nor
evaluates it. The present accepted receiving results are the explicit
complex maps and metrics above. The supplied family's stationary-phase
and boundary-weight providers have not been replaced by that identity.

\begin{thebibliography}{9}
\bibitem{Incoming}
Split-Zero mathematical continuation,
\emph{Complete additive proofs}, 18 September 2026,
sections E, S and J, supplied as \texttt{COMPLETE\_PROOFS.md}.
The corresponding supplied narrative is
\emph{Bidirectional integration of the arithmetic centre, endpoint tensors,
and D032 calculations}, \texttt{INTEGRATION\_REPORT.md}.
\bibitem{Current}
Split-Zero research programme,
\emph{Current cumulative mathematical source},
\texttt{09\_UPDATED\_JOINT\_NOTE.tex}, accepted local successor,
OPG1--5; PQO5; RPC8; PCL1--6.
Exact local paths, byte hashes and extracted line intervals are
recorded in \texttt{SOURCE\_PINS.json}.
\bibitem{Envelopes}
Split-Zero research programme,
\emph{The full tilted mass in the original arithmetic tail}
and \emph{A global tilted-mass transfer with the central source retained},
RTM1--2 and GTM1, preserved provider
\texttt{RTM\_GTM\_original\_providers.tex}.
\bibitem{Deligne}
P. Deligne, \emph{La conjecture de Weil. II},
Publications Math\'ematiques de l'IH\'ES \textbf{52} (1980), 137--252,
\href{https://doi.org/10.1007/BF02684780}{doi:10.1007/BF02684780}.
Full supplied mathematical source retained in the source collection.
\bibitem{DLMF}
R. A. Askey and R. Roy, \emph{Gamma Function}, Chapter 5 of the
NIST Digital Library of Mathematical Functions,
\href{https://dlmf.nist.gov/5.2}{\S5.2(i)}, consulted 19 September 2026.
\end{thebibliography}
\end{document}
